\begin{figure}[ht]
\centering
\begin{tikzpicture}[x=1.0cm, y=1.0cm, font=\small]

% Axes
\draw[->, thick] (0, 0) -- (10.4, 0)
  node[anchor=west, font=\scriptsize] {residual norm $\norm{\mat{A}\vect{x}_{\lambda} - \vect{b}}_{2}$};
\draw[->, thick] (0, 0) -- (0, 6.4)
  node[anchor=south, font=\scriptsize] {solution norm $\norm{\vect{x}_{\lambda}}_{2}$};

% The L-curve: large lambda -> small solution norm, large residual (bottom right);
% small lambda -> large solution norm (noise amplified), small residual (top left)
\draw[very thick, engblue]
  (0.6, 6.0) .. controls (0.9, 3.2) and (1.6, 1.6) ..
  (2.4, 1.15) .. controls (4.5, 0.75) and (7.5, 0.55) .. (9.6, 0.45);

% corner marker (optimal lambda)
\fill[engred] (2.4, 1.15) circle (2pt);
\node[engred, anchor=south west, font=\scriptsize] at (2.5, 1.25) {corner: balanced $\lambda$};

% Over-regularised end
\node[engamber, anchor=west, font=\scriptsize] at (7.0, 0.95)
  {$\lambda$ large: over-smoothed};
\draw[->, engamber, >=stealth] (8.4, 0.85) -- (9.0, 0.55);

% Under-regularised end
\node[enggray, anchor=west, font=\scriptsize] at (1.0, 5.3)
  {$\lambda \to 0$: noise amplified};
\draw[->, enggray, >=stealth] (1.0, 5.6) -- (0.7, 5.9);

\node[anchor=north, font=\scriptsize, text width=11cm, align=center]
  at (5.2, -0.9) {
  Tracing the ridge solution $\vect{x}_{\lambda} =
  (\mat{A}^{\transpose}\mat{A} + \lambda\mat{I})^{-1}
  \mat{A}^{\transpose}\vect{b}$ as $\lambda$ sweeps from large to small
  draws an L. The vertical arm is the ill-conditioned regime where a
  small drop in residual buys a large, noise-driven growth in the
  solution; the horizontal arm is over-smoothing. The corner is the
  defensible choice.
};

\end{tikzpicture}
\caption[The L-curve of Tikhonov regularisation]{The L-curve for a
Tikhonov (ridge) regularised least-squares problem, plotting the
solution norm against the residual norm as the regularisation
parameter $\lambda$ varies. Small $\lambda$ leaves the problem
ill-conditioned and lets measurement noise inflate the solution norm
(the steep vertical arm); large $\lambda$ over-smooths and inflates the
residual (the flat horizontal arm). The corner of the L marks the
parameter that trades the two against each other, and locating it is a
standard, defensible way to choose $\lambda$ without a separate
noise-variance estimate. The regularisation caps the effective
condition number at roughly $\sigma_{1}^{2}/\lambda$, converting a
refusal into a controlled solve.}
\label{fig:vol02:ch09:tikhonov-lcurve}
\end{figure}
